% Chapter introduction: why the topic arose, who first asked, what it answers.
% STATUS: written 2026-10-12; every claim is cited in the sections it points to.

Chapter~\ref{ch:basic-theory} asked what the set of all derivations tells us about
an algebra. This chapter gives the first answer for a class of rings. The class is
the prime rings, which Posner took as the setting of his paper of 1957
\cite[p.~1093]{Posner1957}.

We start with two theorems of Posner. The first says that the composition
of two nonzero derivations of a prime ring is not a derivation
\cite[p.~1093]{Posner1957}. The second says that a nonzero derivation with
$[d(a),a]$ central forces the ring to be commutative
\cite[p.~1093]{Posner1957}. Figure~\ref{fig:prime:posner} shows them.
Section~\ref{sec:prime-semiprime:posner} tells how they were proved and
simplified.

Herstein proved a commutativity theorem for derivations whose values commute
\cite[Theorem~2, p.~369]{Herstein1978}. He also described the elements that
commute with all values of a derivation \cite[p.~509]{Herstein1979}. Section~\ref{sec:prime-semiprime:herstein-school}
follows their commutativity theorems. They need a field that does not lie in
the ring. Martindale built it in 1969 \cite[p.~576]{Martindale1969}, and
Section~\ref{sec:prime-semiprime:quotient-rings} describes it.

With this field, Kharchenko gave a general theory of the identities that
derivations can satisfy in a prime ring \cite[p.~155]{Kharchenko1978}.
Lanski used it to prove the theorems of Posner and their generalizations in a
systematic way \cite[p.~276]{Lanski1988}.
Section~\ref{sec:prime-semiprime:kharchenko} states the theory.

The sections that follow apply it to nilpotent derivations and Engel conditions
(Section~\ref{sec:prime-semiprime:engel-nilpotent}), to Lie ideals
(Section~\ref{sec:prime-semiprime:ideals-lie-ideals}), to rings with involution
(Section~\ref{sec:prime-semiprime:involution}) and to Jordan ideals
(Section~\ref{sec:prime-semiprime:jordan-ideals}).

The literature on this chapter is
large; Appendix~\ref{app:atlas} lists it. Chapter~\ref{ch:jordan-type} and
Chapter~\ref{ch:lie-type} take up the Jordan and Lie variants of the maps studied
here.

\hbfigure{03-prime-semiprime/posner}{The two theorems of Posner \cite[p.~1097]{Posner1957} on derivations of prime rings, the starting point of the chapter.}{fig:prime:posner}
